%---------------------------------------------------------------------
% Part four: the fiqh of salah. Three questions.
%
% Hanafi positions are given first, as the AIOU and Pakistani default.
% Where the imams differ the difference is stated, not hidden -- an
% answer that says "Imam Abu Hanifa held X, the Sahibayn held Y" earns
% more than one that flattens the disagreement away.
%---------------------------------------------------------------------

\section*{حصہ چہارم --- نماز کے مسائل}
\markboth{نماز کے مسائل}{}

%---------------------------------------------------------------------
\begin{sawal}{سوال \textenglish{10} \ \textnormal{\small(خزاں \textenglish{2021}، سوال \textenglish{6})}}
نماز کی شرائط اور ارکان میں کیا فرق ہے؟ نماز کے شرائط، ارکان اور واجبات تفصیل
سے تحریر کریں۔
\end{sawal}

\begin{hal}
\begin{nukta}[پہلے فرق، پھر فہرست --- سوال کا اصل مطالبہ یہی ہے]
\textbf{شرط} نماز سے \emph{باہر} ہوتی ہے اور نماز شروع ہونے سے \emph{پہلے}
پائی جاتی ہے۔ \textbf{رکن} نماز کا \emph{جزو} ہے اور نماز کے \emph{اندر} ادا
ہوتا ہے۔

\smallskip
اس لیے شرط نہ ہو تو نماز \emph{شروع ہی نہیں ہوتی}؛ رکن چھوٹ جائے تو نماز
\emph{باطل ہو جاتی ہے}۔ مثال: وضو شرط ہے، رکوع رکن ہے۔
\end{nukta}

\medskip
\textbf{(الف) شرائط --- چھ:}
\begin{enumerate}\setlength{\itemsep}{1pt}
  \item \textbf{طہارت} --- بدن، کپڑے اور جگہ تینوں کی پاکی۔
  \item \textbf{سترِ عورت} --- مرد کے لیے ناف سے گھٹنے تک؛ عورت کے لیے چہرے،
        ہاتھوں اور قدموں کے سوا سارا بدن۔
  \item \textbf{استقبالِ قبلہ}۔
  \item \textbf{وقت} --- ہر نماز اپنے مقررہ وقت میں۔
  \item \textbf{نیت} --- دل کا ارادہ؛ زبان سے کہنا شرط نہیں۔
  \item \textbf{تکبیرِ تحریمہ} --- حنفیہ کے نزدیک یہ \emph{شرط} ہے (کیونکہ یہ
        نماز میں داخل ہونے کی کنجی ہے)، جمہور کے نزدیک \emph{رکن}۔ یہ اختلاف
        لکھ دینا جواب کو مضبوط کرتا ہے۔
\end{enumerate}

\medskip
\textbf{(ب) ارکان / فرائض --- چھ:}
\begin{enumerate}\setlength{\itemsep}{1pt}
  \item \textbf{قیام} --- فرض قراءت کے بقدر کھڑا ہونا۔
  \item \textbf{قراءت} --- کم از کم ایک آیت۔
  \item \textbf{رکوع}۔
  \item \textbf{سجود} --- ہر رکعت میں دو۔
  \item \textbf{قعدۂ اخیرہ} --- تشہد پڑھنے کے بقدر بیٹھنا۔
  \item \textbf{خروج بصنعہ} --- اپنے فعل سے نماز سے نکلنا، یعنی سلام۔
\end{enumerate}

\medskip
\textbf{(ج) واجبات --- اہم دس:} سورۃ فاتحہ پڑھنا \ \textbullet\ \ فاتحہ کے بعد
سورت یا تین آیات ملانا \ \textbullet\ \ فاتحہ کو سورت پر مقدم رکھنا \
\textbullet\ \ قومہ (رکوع سے سیدھا کھڑا ہونا) \ \textbullet\ \ جلسہ (دو سجدوں
کے درمیان بیٹھنا) \ \textbullet\ \ تعدیلِ ارکان (ہر رکن میں کم از کم ایک تسبیح
کا وقفہ) \ \textbullet\ \ قعدۂ اولیٰ \ \textbullet\ \ دونوں قعدوں میں تشہد \
\textbullet\ \ لفظ ''السلام'' سے نماز ختم کرنا \ \textbullet\ \ وتر میں دعائے
قنوت۔

\medskip
\begin{ghalti}[یہ فرق نمبروں کا فیصلہ کرتا ہے]
\textbf{فرض} چھوٹ جائے تو سجدہ سہو کافی نہیں --- نماز \emph{دوبارہ} پڑھنی پڑے
گی۔ \textbf{واجب} بھول کر چھوٹے تو \emph{سجدہ سہو} سے تلافی ہو جاتی ہے۔ اور
\textbf{جان بوجھ کر} واجب چھوڑا جائے تو سجدہ سہو کافی نہیں، اعادہ واجب ہے۔
\end{ghalti}
\end{hal}

%---------------------------------------------------------------------
\begin{sawal}{سوال \textenglish{11} \ \textnormal{\small(بہار \textenglish{2013}، سوال \textenglish{6} اور سوال \textenglish{8})}}
''اوقات الصلوٰۃ'' کے مسائل پر تفصیلی نوٹ لکھیں، نیز مکروہاتِ نماز بیان کریں۔
\end{sawal}

\begin{hal}
\textbf{(الف) پانچ نمازوں کے اوقات۔}

\begin{center}
\begin{tabular}{@{}p{2.4cm} p{12.0cm}@{}}
\toprule
\textbf{فجر}   & صبح صادق (مشرق میں سفیدی) سے طلوعِ آفتاب تک۔ \\
\textbf{ظہر}   & زوالِ آفتاب سے۔ آخری وقت میں اختلاف ہے: امام ابو حنیفہؒ کے
                 نزدیک جب سایہ \emph{دو مثل} ہو جائے، صاحبین و جمہور کے نزدیک
                 \emph{ایک مثل} پر۔ \\
\textbf{عصر}   & ظہر کا وقت ختم ہونے سے غروبِ آفتاب تک۔ \\
\textbf{مغرب}  & غروبِ آفتاب سے شفق غائب ہونے تک۔ \\
\textbf{عشاء}  & شفق غائب ہونے سے صبح صادق تک۔ (شفق سفیدی ہے یا سرخی --- دیکھیے
                 حصہ اول۔) \\
\textbf{وتر}   & عشاء کے بعد صبح صادق تک۔ \\
\bottomrule
\end{tabular}
\end{center}

\medskip
\textbf{(ب) تین اوقاتِ مکروہہ --- ان میں کوئی نماز جائز نہیں:} حضرت عقبہ بن
عامرؓ فرماتے ہیں کہ نبی کریم \saw نے ہمیں تین وقتوں میں نماز پڑھنے سے منع فرمایا
(مسلم):

\begin{enumerate}\setlength{\itemsep}{1pt}
  \item طلوعِ آفتاب کے وقت --- جب تک سورج بلند نہ ہو جائے۔
  \item عین نصف النہار (استواء) --- جب سورج ٹھیک سر پر ہو۔
  \item غروبِ آفتاب کے وقت --- جب سورج ڈوب رہا ہو۔
\end{enumerate}

\smallskip
\textbf{ایک استثنا:} اگر عصر کی فرض نماز رہ گئی ہو تو غروب کے وقت بھی پڑھی جا
سکتی ہے۔

\medskip
\textbf{(ج) دو اوقات جن میں صرف نفل مکروہ ہے:} فجر کے بعد طلوعِ آفتاب تک، اور
عصر کے بعد غروب تک۔ فرض، قضا اور جنازہ ان میں جائز ہیں (بخاری و مسلم)۔

\medskip
\textbf{(د) مکروہاتِ نماز۔} یہ وہ کام ہیں جن سے نماز باطل نہیں ہوتی، مگر ثواب
کم ہو جاتا ہے:

\begin{itemize}\setlength{\itemsep}{1pt}
  \item \textbf{عبث حرکت اور التفات} --- کپڑے سے کھیلنا، ادھر ادھر دیکھنا:
        ''وہ چوری ہے جو شیطان نماز سے کرتا ہے'' (بخاری)۔
  \item \textbf{کفِ ثوب و سدل، اقعاء} --- آستین سمیٹنا، چادر بغیر لپیٹے
        لٹکانا، کتے کی طرح بیٹھنا (بخاری)۔
  \item بلا عذر ٹیک لگا کر پڑھنا؛ \textbf{کھانا سامنے یا حاجت کا تقاضا ہو:}
        ''\ar{لَا صَلَاةَ بِحَضْرَةِ الطَّعَامِ}'' (مسلم)۔
\end{itemize}
\end{hal}

%---------------------------------------------------------------------
\begin{sawal}{سوال \textenglish{12} \ \textnormal{\small(بہار \textenglish{2012}، سوال \textenglish{4} اور \textenglish{8}؛ بہار \textenglish{2013}، سوال \textenglish{8})}}
نماز جمعہ، صلوٰۃ المریض، قیامِ شہر رمضان اور نماز استسقاء کے مسائل بیان کریں،
نیز سجدہ سہو واجب ہونے کی پانچ صورتیں لکھیں۔
\end{sawal}

\begin{hal}
\textbf{(الف) نماز جمعہ۔} حکم قرآن میں ہے: \ar{يَا أَيُّهَا الَّذِينَ آمَنُوا
إِذَا نُودِيَ لِلصَّلَاةِ مِن يَوْمِ الْجُمُعَةِ فَاسْعَوْا إِلَىٰ ذِكْرِ
اللهِ وَذَرُوا الْبَيْعَ} (الجمعہ، \textenglish{9})۔

\smallskip
\textbf{شرائطِ وجوب} (کس پر فرض ہے): مرد ہونا، آزاد، مقیم (مسافر نہ ہو)، تندرست،
بینا اور چلنے کی قدرت رکھنے والا۔ چنانچہ عورت، مسافر، مریض، نابینا اور معذور پر
فرض نہیں --- البتہ پڑھ لیں تو ادا ہو جائے گی۔

\smallskip
\textbf{شرائطِ صحت} (نماز درست ہونے کی): شہر یا قصبہ ہونا، وقتِ ظہر، \textbf{خطبہ}
(یہ شرط ہے)، جماعت، اور \textbf{اذنِ عام} --- یعنی دروازے سب کے لیے کھلے ہوں۔ فرض
دو رکعت ہیں۔

\medskip
\textbf{(ب) صلوٰۃ المریض۔} اصول قرآن میں ہے: \ar{فَاتَّقُوا اللهَ مَا
اسْتَطَعْتُمْ} (التغابن، \textenglish{16})، اور ترتیب حضرت عمران بن حصینؓ کی
حدیث میں: ''\ar{صَلِّ قَائِمًا، فَإِنْ لَمْ تَسْتَطِعْ فَقَاعِدًا، فَإِنْ لَمْ
تَسْتَطِعْ فَعَلَىٰ جَنْبٍ}'' (بخاری)۔

\begin{itemize}\setlength{\itemsep}{1pt}
  \item کھڑا نہ ہو سکے تو \textbf{بیٹھ کر}؛ بیٹھ نہ سکے تو \textbf{لیٹ کر
        اشارے سے} (چت لیٹ کر پاؤں قبلہ کی طرف، یا کروٹ پر)۔
  \item رکوع و سجدہ نہ کر سکے تو سر کے اشارے سے --- سجدے کا اشارہ رکوع سے
        \emph{نیچا} ہو۔
  \item اشارہ بھی ممکن نہ ہو اور یہ کیفیت چھ نمازوں سے زیادہ رہے تو قضا ساقط۔
\end{itemize}

\medskip
\textbf{(ج) قیامِ شہر رمضان (تراویح)۔} فضیلت: ''\ar{مَنْ قَامَ رَمَضَانَ
إِيمَانًا وَاحْتِسَابًا غُفِرَ لَهُ مَا تَقَدَّمَ مِنْ ذَنْبِهِ}'' (بخاری و
مسلم)۔ مرد و عورت دونوں کے لیے \textbf{سنتِ مؤکدہ}، حنفیہ کے نزدیک بیس رکعت۔
نبی کریم \saw نے تین رات باجماعت پڑھائی، پھر فرض سمجھے جانے کے اندیشے سے چھوڑ
دی۔ حضرت عمرؓ نے مستقل اجتماعی صورت دی: ''\ar{نِعْمَتِ الْبِدْعَةُ هَذِهِ}''
(بخاری)۔

\medskip
\textbf{(د) نماز استسقاء (بارش کی دعا)۔} \textbf{ائمہ کا اختلاف:} امام ابو
حنیفہؒ کے نزدیک صرف \textbf{دعا و استغفار}، باجماعت نماز سنت نہیں؛ صاحبین اور
جمہور کے نزدیک \textbf{دو رکعت باجماعت} سنت ہے۔ طریقہ: تین دن روزہ و توبہ، پھر
کھلے میدان میں دو رکعت، خطبہ اور \textbf{تحویلِ رداء} (چادر پلٹنا)۔ بنیاد:
\ar{اسْتَغْفِرُوا رَبَّكُمْ إِنَّهُ كَانَ غَفَّارًا يُرْسِلِ السَّمَاءَ
عَلَيْكُم مِّدْرَارًا} (نوح، \textenglish{10--11})۔

\medskip
\textbf{(ہ) سجدہ سہو واجب ہونے کی پانچ صورتیں۔} یہ سوال بہار \textenglish{2012}
میں انہی الفاظ میں آیا:

\begin{enumerate}\setlength{\itemsep}{1pt}
  \item کوئی \textbf{واجب بھول کر چھوٹ جائے} --- مثلاً قعدۂ اولیٰ یا تشہد۔
  \item کسی واجب میں \textbf{تاخیر} ہو جائے --- مثلاً رکوع یا سجدے میں غیر
        معمولی توقف۔
  \item کسی رکن کا \textbf{تکرار} ہو جائے --- ایک رکعت میں دو رکوع یا تین سجدے۔
  \item ارکان کی \textbf{ترتیب بدل} جائے --- مثلاً سجدے سے پہلے قعدہ کر لینا۔
  \item قراءت کی \textbf{کیفیت بدل} جائے --- جہری نماز میں آہستہ یا سری میں
        بلند پڑھنا۔
\end{enumerate}

\smallskip
\textbf{طریقہ:} آخری قعدے میں تشہد کے بعد ایک طرف سلام پھیر کر دو سجدے کیجیے،
پھر تشہد، درود اور دعا پڑھ کر سلام پھیر دیجیے۔
\end{hal}
